% 09-evaluation.tex
\section{Evaluation: Readability, Correctness, and Overhead}
\label{sec:eval}

This chapter evaluates TensorPlay on three axes that a readable framework must earn: how small the core is to read, how precisely it guards correctness when the training loop mutates tensors, and how little it costs when every guard and profiler probe is disabled. It also adds four evidence blocks that complete the architecture story: the populated CUDA/Vulkan backend slots (\S\ref{subsec:eval-backends}), the support cores that sit below the kernels (\S\ref{subsec:eval-supportcores}), the three-track test suite (\S\ref{subsec:eval-tests}), and measured benchmark numbers (\S\ref{subsec:eval-bench}).

%======================================================================
\subsection{Readability (Lines of Core)}
\label{subsec:eval-readability}

Table~\ref{tab:evalsize} quantifies the ``bone vs.\ flesh'' claim of \S\ref{sec:arch}: the architectural skeleton---object model, dispatcher layering, \texttt{TensorIterator}, autograd \texttt{GraphTask} engine---is retained line-for-line, while production axes (op coverage, dispatch depth, dtype$\times$op completeness, sparse/quantized) are trimmed.

\begin{table}[H]
\centering
\caption{Core size. TensorPlay intentionally trims production axes. ``Core'' is the four C++ pillars plus bindings; generated code and tests are excluded.}
\label{tab:evalsize}
\small
\begin{tabular}{lrr}
\toprule
Metric & TensorPlay & Broader implementation \\
\midrule
C++ core & $\approx$ 75\,k & hundreds of k (core runtime $>$ 300\,k) \\
Python surface & 267\,k / 899 files (see Table~\ref{tab:pybreakdown}) & larger package surface \\
Operators (unique \texttt{m.impl}) & 642--673 & 2000+ \\
\texttt{DispatchKey} cardinality & 13 & 60+ \\
Autograd nodes (generated) & 148 & thousands \\
Build artifacts & 4 libs: \texttt{libp10}, \texttt{libtpx}, \texttt{libstax}, \texttt{\_C} & dozens \\
\bottomrule
\end{tabular}
\end{table}

\begin{table}[H]
\centering
\caption{Where the 75\,k C++ lines live. Generated files are excluded.}
\label{tab:corebreakdown}
\small
\begin{tabular}{llr}
\toprule
Component & Representative content & Approx.\ lines \\
\midrule
\texttt{p10} object model & \texttt{Tensor}, \texttt{TensorImpl}, \texttt{Storage}, version counter & $\sim$18\,k \\
\texttt{p10} dispatcher & 13-slot table, key set, registration macro & $\sim$1.5\,k \\
\texttt{p10} kernels & CPU backends ($\ast$.cpp) with CUDA counterparts ($\ast$.cu) & $\sim$35\,k \\
\texttt{tpx} autograd engine & \texttt{Node}, \texttt{Edge}, \texttt{GraphTask}, \texttt{ReadyQueue}, \texttt{Engine} & $\sim$6\,k \\
\texttt{stax} + compiler & graph IR, fusion, interpreter & $\sim$8\,k \\
Bindings & CPython bridge, per-domain units & $\sim$7\,k \\
\bottomrule
\end{tabular}
\end{table}

The size gap is not a virtue in itself; it is the budget that buys readability. Every indirection that remains is load-bearing and has a single definition site.

%======================================================================
\subsection{Where the Python Lines Live}
\label{subsec:eval-pybreakdown}

Table~\ref{tab:evalsize} cites a single Python number, which invites the question ``what is inside it?''. Table~\ref{tab:pybreakdown} decomposes the 267\,k lines so the ``thin veneer'' claim can be inspected subsystem by subsystem.

\begin{table}[H]
\centering
\caption{Python surface decomposition. ``Model zoos'' are pretrained-weight wrappers for third-party checkpoints; they exercise the \texttt{\_C} surface but add no framework mechanism.}
\label{tab:pybreakdown}
\small
\begin{tabular}{llrr}
\toprule
Subpackage & What lives there & Files & Lines \\
\midrule
\texttt{distributed/} & c10d-style collectives, fsdp, DTensor, elastic, rpc & 367 & 62{,}770 \\
\texttt{vision/} + \texttt{audio/} & datasets, transforms, model zoos & 148 & 51{,}039 \\
\texttt{graph/} & FX-style symbolic tracer, Node/Graph, symbolic shapes & 113 & 27{,}991 \\
\texttt{nn/} & modules, functional, init & 42 & 22{,}832 \\
\texttt{\_stax/} & shape prop, DCE, const fold, interpreter & 20 & 11{,}528 \\
\texttt{optim/} & 15 optimizers, foreach/fused paths & 24 & 11{,}365 \\
\texttt{export/} + \texttt{onnx/} & draft-export, effect-token passes & 34 & 9{,}216 \\
\texttt{autograd/} + \texttt{\_transforms/} & Function, grad-mode, vmap wrappers & 24 & 8{,}072 \\
Other 17 subpackages & cuda, profiler, amp, sparse, quantization, \dots & 134 & 35{,}884 \\
Root modules & \texttt{\_tensor.py}, \texttt{library.py} (1\,779), \texttt{overrides.py} & 23 & 37{,}805 \\
\midrule
Total & & 899 & 266{,}907 \\
\bottomrule
\end{tabular}
\end{table}

Two observations follow. First, the framework-mechanism content---\texttt{nn}, \texttt{optim}, \texttt{\_stax}, \texttt{autograd}, root modules---is $\sim$91\,k lines (34\%), in line with the ``thin veneer'' claim of \S\ref{sec:arch}: the tensor type and the library-impl registration surface are the only large hand-written integration points. Second, the single largest block is \texttt{distributed/} (63\,k, 24\%); its FSDP subtree alone is 9.5\,k lines. The model zoos (\texttt{vision/} + \texttt{audio/}, 51\,k) are dataset/checkpoint plumbing that could be deleted without touching dispatch or autograd; they exist so the end-to-end benchmarks of \S\ref{subsec:eval-bench} have realistic workloads. None of the 899 files registers a kernel or a dispatch key; everything funnels back into \texttt{\_C}.

%======================================================================
\subsection{GPU Backends: CUDA and Vulkan Slots}
\label{subsec:eval-backends}

\S\ref{sec:redispatch} declares 13 keys with \texttt{CUDA}=1 and \texttt{Vulkan}=2 as dense-backend slots, and \S\ref{subsec:eval-microbench} cites their dispatch cost. This subsection shows that the slots are populated: backends exist, not just keys.

\begin{table}[H]
\centering
\caption{Backend inventory. ``Registrations'' counts \texttt{TENSORPLAY\_LIBRARY\_IMPL} macro expansions; ``unique ops'' deduplicates \texttt{m.impl} names per backend subtree.}
\label{tab:backends}
\small
\begin{tabular}{lrrrr}
\toprule
Backend & Files & Lines & Registrations & Unique ops \\
\midrule
CPU & 69 & 61{,}108 & 66 & 1{,}325 \\
CUDA & 86 & 56{,}949 & 66 & 1{,}274 \\
Vulkan & 93 & 16{,}289 + 4{,}500 GLSL & 26 & 145 \\
Composite & 34 & 8{,}533 & 34 & --- \\
\bottomrule
\end{tabular}
\end{table}

\begin{figure}[H]
\centering
\begin{tikzpicture}[every node/.style={font=\scriptsize}]
% unique m.impl ops per backend
\def\cw{0.62}
\def\ch{0.0028}
\foreach \x/\name/\count/\fillc/\drawc in {
  0/CPU/1325/blue!8/tpblue,
  1/CUDA/1274/orange!10/tporange,
  2/Vulkan/145/purple!8/tppurple}{
  \pgfmathsetmacro{\h}{\count*\ch}
  \pgfmathsetmacro{\lbl}{\h+0.55}
  \fill[\fillc] (\x*\cw+0.05,0) rectangle (\x*\cw+0.05+0.5,\h);
  \draw[\drawc] (\x*\cw+0.05,0) rectangle (\x*\cw+0.05+0.5,\h);
  \node at (\x*\cw+0.3,\lbl) {\name: \count};
  \node at (\x*\cw+0.3,-0.25) {\texttt{m.impl}};
}
% axis
\draw[->,gray] (-0.15,0) -- (2.35,0);
\node[anchor=west,gray] at (2.45,3.7) {unique ops};
\end{tikzpicture}
\caption{Unique op names per backend (Table~\ref{tab:backends}). The CUDA column is 96\% of the CPU column; Vulkan is the deliberately small teaching backend.}
\label{fig:backendops}
\end{figure}

\subsubsection{CUDA backend (56.9\,k lines, 1{,}274 ops)}

The CUDA backend keeps a one-file-per-family layout parallel to the CPU tree where kernels are separable: arithmetic, pointwise (2{,}154 lines), reduction, convolution (2{,}477), indexing (3{,}362), plus CUDA-only infrastructure: a stream-aware caching allocator (1{,}288 lines), stream/device bookkeeping, a cuBLAS handle pool, NCCL communicator cache (418 lines), a Philox-based generator, and CUDA-graph capture. Every kernel file ends with one registration block, e.g.\ \texttt{TENSORPLAY\_LIBRARY\_IMPL(CUDA, PointwiseKernels)}, so each translation unit is independently iterable.

\paragraph{Tiered reduction engine.} Reductions (sum, mean, min/max, Welford variance) share one code-generation template with a three-tier combine strategy: intra-warp via \texttt{\_\_shfl\_down\_sync} (no shared memory), inter-warp via shared memory within the block, and a global second-pass kernel when the tensor spans many blocks. A traits type selects the accumulation dtype per input dtype---\texttt{float} accumulation for \texttt{Half}/\texttt{BFloat16} inputs keeps precision---and per-operation functors (\texttt{SumOps}, \texttt{MeanOps}, \texttt{MinMaxOps}, \texttt{WelfordOps}) plug into the same launcher.

\begin{lstlisting}[caption={The intra-warp tier of the shared reduction template; complex numbers shuffle their two lanes because the shuffle intrinsic has no complex overload.}, label=lst:cudareduce]
__device__ inline float shfl_down(float value, uint32_t mask, int offset) {
  return __shfl_down_sync(mask, value, offset);
}
// complex has no intrinsic __shfl_down_sync overload; shuffle the
// real and imaginary lanes separately.
__device__ inline complex_t shfl_down(complex_t value,
                                       uint32_t mask, int offset) {
  float re = __shfl_down_sync(mask, value.real(), offset);
  float im = __shfl_down_sync(mask, value.imag(), offset);
  return complex_t(re, im);
}
\end{lstlisting}

\paragraph{Multi-tensor apply for optimizers.} Optimizer steps act on hundreds of parameter tensors of the same dtype. The \emph{foreach} path chunks work from the whole tensor list into one grid-stride kernel (\texttt{foreach\_mta::launch}), so an Adam step updates weights, moments, and buffers in a handful of grouped launches instead of hundreds of individual ones. Eligibility is structural---same device, same dtype---and ineligible groups fall back to per-tensor launches. The CPU benchmark of \S\ref{subsec:eval-bench} (optimizer geomean 3.83$\times$) is largely attributable to this path plus fused one-kernel steps.

\paragraph{cuSOLVER decompositions.} Beyond GEMM (cuBLAS, column-major staging via \texttt{clone\_batched\_column\_major}), the backend wraps cuSOLVER for SVD, Cholesky, and LU with device-side info pointers so positive-definiteness failures surface as exceptions rather than silent garbage.

CUDA-graph capture shows the backend's care for ordering. Capture must begin on a non-default stream, one capture per thread, and pool routing must be armed \emph{before} the driver begins recording:

\begin{lstlisting}[caption={CUDA-graph capture entry: nesting and reuse are rejected loudly, and the side stream is chosen before any driver call.}, label=lst:cudagraph]
void CUDAGraph::capture_begin(uint64_t pool_id, CaptureMode mode,
                              const CUDAStream& stream) {
    GraphState& state = GraphState::instance();
    const int device = currentDevice();
    CUDAStream side = CUDAStream::undefined();
    {
        std::lock_guard<std::mutex> lock(state.mutex);
        if (state.capturing.count(std::this_thread::get_id())) {
            TP_THROW(RuntimeError,
                     "nested CUDA graph capture is not supported on one "
                     "thread (captures on different devices/threads may "
                     "run concurrently)");
        }
        if (exec_ != nullptr || graph_ != nullptr || pool_id_ != 0 ||
            owns_pool_ref_) {
            TP_THROW(RuntimeError,
                     "this CUDAGraph already holds a capture or executable; "
                     "create a new one");
        }
        /* ... arm memory-pool routing, then cudaStreamBeginCapture ... */
\end{lstlisting}

The RNG half of the story is capture-awareness. Kernels draw randoms from a Philox counter, and during graph capture the counter cannot be baked into the kernel launch---each replay must see a fresh slice of the random stream. The generator therefore exposes two regimes:

\begin{lstlisting}[caption={Philox counter reservation, in two regimes: eager launches get plain (seed, offset); captured kernels read them from device buffers that a replay prologue refreshes.}, label=lst:philox]
// Atomically reserves `increment` philox counter values for the current
// device and returns the (seed, offset) the launching kernel should consume.
P10_API std::pair<uint64_t, uint64_t> philox_engine_inputs(uint64_t increment);

// graph capture kernels consume the plain (seed, offset) values; while a
// capture is underway they instead read seed/offset from device buffers owned
// by the capturing graph plus a per-kernel intragraph offset, so replay
// prologues can refresh the buffers and every replay consumes a fresh slice
// of the random stream instead of repeating capture-time values.
P10_API PhiloxCudaState philox_cuda_state(uint64_t increment);
\end{lstlisting}

% ---- Figure: Philox RNG kernel flow (CUDA) ----
\begin{figure}[H]
\centering
\begin{tikzpicture}[every node/.style={font=\scriptsize\sffamily}]
\node[libbox, fill=blue!8, draw=tpblue, minimum width=3.1cm, minimum height=0.75cm] (state) {PhiloxCudaState (seed, offset)};
\node[libbox, fill=blue!8, draw=tpblue, below=0.5cm of state, minimum width=3.1cm, minimum height=0.75cm] (kernel) {elementwise grid-stride kernel};
\node[libbox, fill=orange!10, draw=tporange, below=0.5cm of kernel, minimum width=3.1cm, minimum height=0.75cm] (init) {\texttt{curand\_init}};
\node[libbox, fill=orange!10, draw=tporange, below=0.5cm of init, minimum width=3.1cm, minimum height=0.75cm] (loop) {grid-stride loop};
\node[libbox, fill=green!8, draw=tpgreen, below=0.5cm of loop, minimum width=3.1cm, minimum height=0.75cm] (dist) {distribution functors};
\node[draw=none, right=0.7cm of dist, text width=3.4cm, align=left, font=\tiny] {e.g.\ \texttt{curand\_uniform4} / \texttt{curand\_normal4}: four counter values per call, vectorized};
\draw[arrow, color=tpblue] (state) -- (kernel);
\draw[arrow, color=tporange] (kernel) -- (init);
\draw[arrow, color=tporange] (init) -- (loop);
\draw[arrow, color=tpgreen] (loop) -- (dist);
\end{tikzpicture}
\caption{The random-sampling path: each launch unpacks the capture-aware counter state, seeds \texttt{curand}, and draws vectorized samples inside one grid-stride kernel.}
\label{fig:philoxflow}
\end{figure}



\subsubsection{Vulkan backend (16.3\,k C++ + 4.5\,k GLSL, 145 ops)}

Vulkan is the smallest complete backend, which is exactly its teaching value: it demonstrates that slot 2 of the three backend keys accepts a new device family with \emph{no dispatcher change}. The layout is three layers:

\begin{itemize}[leftmargin=1.2em]
\item An \texttt{api} layer (29 files): a thin object model over raw handles---\texttt{Adapter} (device selection), \texttt{Context}, \texttt{Resource} (buffer/image with VMA), \texttt{Pipeline}, \texttt{Command}, \texttt{QueryPool}, \texttt{vTensor} (the backend's tensor view).
\item An \texttt{ops} layer (34 files): 145 \texttt{m.impl} registrations, one closing \texttt{TENSORPLAY\_LIBRARY\_IMPL(Vulkan,...)} block per file (26 blocks total). Kernels fetch shaders by name through a registry lookup \texttt{VK\_KERNEL(name)}.
\item A \texttt{glsl} layer (70 template files): compute shaders with a mini preprocessor---\texttt{\$\{PARAM\}} substitution and \texttt{\$if}/\texttt{\$elif}/\texttt{\$else} blocks---driven by 63 variant-table entries.
\end{itemize}

The build pipeline is fully generated: a 418-line Python stage preprocesses each template, invokes \texttt{glslangValidator} for Vulkan 1.0 SPIR-V, parses the descriptor-set bindings, and emits an embedded C++ registry; the build rule wires the dependency so editing any shader or variant table re-embeds all 220 compiled \texttt{.spv} variants. The quantized family is a first-class citizen here (quantized convolution, per-tensor dequantize, \dots), which is why Vulkan ships quantized kernels while CPU keeps quantization the Python FakeQuantize layer (\S\ref{subsec:eval-trimmed}).

The matmul kernel shows the whole layering in one body---validate, allocate a device tensor, pack layouts, bind uniforms, submit:

\begin{lstlisting}[caption={Vulkan matmul kernel body: the three-layer pattern (validate $\rightarrow$ \texttt{vTensor} + packing $\rightarrow$ uniform block $\rightarrow$ submit).}, label=lst:vulkanmm]
Tensor mm_kernel(const Tensor& self, const Tensor& mat2) {
  TP_CHECK(self.dtype() == DType::Float32 && mat2.dtype() == DType::Float32,
          "Vulkan mm supports Float32 tensors only");
  TP_CHECK(self.dim() == 2, "Vulkan mm: self must be a matrix");
  TP_CHECK(mat2.dim() == 2, "Vulkan mm: mat2 must be a matrix");
  TP_CHECK(self.size(1) == mat2.size(0),
           "Vulkan mm: matrix dimensions do not match");

  api::Context* const context = api::context();
  const int64_t M = self.size(0);
  const int64_t N = mat2.size(1);
  const int64_t K = self.size(1);
  api::vTensor v_output{context, {M, N}, self.dtype()};
  /* pack channels/height layouts, fill the MMBlock uniform
     (out_sizes, step_size = ceil(K/4)), and submit a 4x4-tile grid */
\end{lstlisting}

%======================================================================
\subsection{AMD Portability: the Hipify Build Path}
\label{subsec:eval-rocm}

Beyond the three first-class backends, TensorPlay runs on AMD GPUs through an automated translation of its CUDA sources rather than a maintained fork. At configuration time a build script stages every CUDA source and binding into a separate tree and applies regular-expression rewrites---\texttt{cudaMalloc} becomes \texttt{hipMalloc}, kernel launches become \texttt{hipLaunchKernelGGL}---after which the staged tree compiles with the HIP toolchain. Two build options enforce mutual exclusion between the CUDA and ROCm paths because both register the same device and dispatch slots; the dispatcher then routes to the translated kernels transparently. The design lesson: when a backend differs only in vocabulary, translate the vocabulary at build time and keep one semantic source, rather than growing a hand-maintained sibling tree that drifts.

%======================================================================
\subsection{Numerical Stability: Cascade Sums on CPU}
\label{subsec:eval-cascadesum}

Long float sums accumulate rounding error linearly in the element count when a single accumulator is used. The CPU reduction path implements a \emph{cascade} scheme instead: partial sums are kept in a small array of levels, and a level is flushed into the next level whenever its index hits a power-of-two step. Error growth drops from $O(n)$ to $O(\log n)$ because each level only ever adds numbers of comparable magnitude.

\begin{lstlisting}[caption={The cascade accumulator: four levels, flushed on a power-of-two cadence; vector loads feed every lane in parallel.}, label=lst:cascadesum]
constexpr int64_t levels = 4;
const int64_t level_power = std::max<int64_t>(4, ceil_log2(size) / levels);
const int64_t level_step  = int64_t(1) << level_power;
const int64_t level_mask  = level_step - 1;
std::array<std::array<T, N>, levels> acc{};   // one accumulator per lane
for (; i + level_step <= size;) {
  /* load N-wide vector, add into level 0 */
  if ((i & level_mask) == 0) {                // cadence hit: flush upward
    for (int64_t l = 0; l + 1 < levels; ++l) {
      /* acc[l+1] += acc[l]; acc[l] = 0; */
    }
  }
}
\end{lstlisting}

The AVX-512 fast path for \texttt{normsq} (sum of squares) fuses multiply and accumulate with FMA instructions; on compatible hardware this is the mechanism behind the $9.9\times$ \texttt{norm2} speedup recorded in Figure~\ref{fig:benchhist}, where the reference's scalar loop cannot use the same instruction mix.

%======================================================================
\subsection{Distributed Training: FSDP, DTensor, Functional Collectives}
\label{subsec:eval-distributed}

The distributed stack (63\,k Python lines plus 17\,k C++ lines of process-group backends over gloo and MPI) is organized in three layers that mirror the pillar philosophy: a C++ communication layer, a differentiable Python functional layer, and composable sharding policies.

\paragraph{FSDP.} Fully sharded data parallelism is configured by a small enum vocabulary rather than ad-hoc flags: \texttt{FULL\_SHARD} (parameters sharded, gathered per layer, freed after use), \texttt{SHARD\_GRAD\_OP} (gather once per forward, shard only gradients and optimizer state), \texttt{NO\_SHARD} (plain DDP), and a hybrid intra/inter-node variant. Policies compose as dataclasses---\texttt{MixedPrecision} chooses independent dtypes for parameters, reduction, and buffers; \texttt{CPUOffload} swaps parameters to host memory between uses; \texttt{BackwardPrefetch} overlaps the next layer's gather with the current layer's backward.

\begin{lstlisting}[caption={The strategy and policy vocabulary of FSDP: everything the sharding engine needs is one enum plus three small dataclasses.}, label=lst:fsdpvocab]
class ShardingStrategy(Enum):
    FULL_SHARD    = auto()   # params sharded, gathered per-layer
    SHARD_GRAD_OP = auto()  # params gathered once; grads+state sharded
    NO_SHARD      = auto()   # plain DDP
    HYBRID_SHARD  = auto()   # replicate intra-node, shard inter-node

@dataclass
class MixedPrecision:
    param_dtype: Any = None     # computation dtype
    reduce_dtype: Any = None    # all-reduce dtype
    buffer_dtype: Any = None    # communication buffers

@dataclass
class CPUOffload:
    offload_params: bool = False
\end{lstlisting}

\paragraph{DTensor.} Sharded tensors are described by a \emph{placement} algebra over a device mesh: \texttt{Shard(dim)} splits along a dimension, \texttt{Replicate} copies, and \texttt{Partial(reduce\_op)} holds partial values that reduce to the global tensor. Operations on a \texttt{DTensor} dispatch through a sharding-aware layer that inserts the minimal set of collectives (\texttt{all\_gather}, \texttt{reduce\_scatter}) needed to satisfy the output placement---the same redispatch idea as \S\ref{sec:redispatch}, lifted from backends to topologies.


% Figure: functional collective across the Python/C++ boundary
\begin{figure}[H]
\centering
\begin{tikzpicture}[every node/.style={font=\scriptsize\sffamily}]
\node[libbox, fill=blue!8, draw=tpblue, minimum width=4.6cm, minimum height=0.7cm] (py) {Python: \texttt{all\_reduce(t, op="sum")}};
\node[libbox, fill=orange!8, draw=tporange, below=0.55cm of py, minimum width=4.6cm, minimum height=0.7cm] (bind) {C++ bindings};
\node[libbox, fill=orange!8, draw=tporange, below=0.55cm of bind, minimum width=4.6cm, minimum height=0.7cm] (pg) {process group (gloo / MPI / NCCL)};
\node[libbox, fill=gray!12, draw=gray, below=0.55cm of pg, minimum width=4.6cm, minimum height=0.7cm] (work) {asynchronous work handle};
\node[libbox, fill=green!8, draw=tpgreen, below=0.55cm of work, minimum width=4.6cm, minimum height=0.7cm] (async) {\texttt{AsyncCollectiveTensor}};
\draw[arrow, color=tpblue] (py) -- node[right, font=\tiny]{tensor list} (bind);
\draw[arrow, color=tporange] (bind) -- (pg);
\draw[arrow, color=gray] (pg) -- (work);
\draw[arrow, color=tpgreen] (work) -- node[right, font=\tiny]{waits lazily} (async);
\end{tikzpicture}
\caption{One collective across the language boundary: the handle returns immediately; the wrapper demotes itself to a plain tensor at the first real use.}
\label{fig:collbridge}
\end{figure}

% Figure: DTensor sharding resolution
\begin{figure}[H]
\centering
\begin{tikzpicture}[every node/.style={font=\scriptsize\sffamily}, node distance=0.45cm]
\node[libbox, fill=blue!8, draw=tpblue, minimum width=4.4cm, minimum height=0.75cm] (user) {\texttt{matmul(A, B)} on DTensors};
\node[libbox, fill=purple!8, draw=tppurple, below=of user, minimum width=4.4cm, minimum height=0.75cm] (check) {check \texttt{A.placement}, \texttt{B.placement}};
\node[libbox, fill=purple!8, draw=tppurple, below=of check, minimum width=4.4cm, minimum height=0.75cm] (reshard) {insert collectives to reshard};
\node[libbox, fill=green!8, draw=tpgreen, below=of reshard, minimum width=4.4cm, minimum height=0.75cm] (local) {local op on shards (native kernels)};
\node[libbox, fill=blue!8, draw=tpblue, below=of local, minimum width=4.4cm, minimum height=0.75cm] (out) {output DTensor with resolved placement};
\node[draw=none, right=0.55cm of reshard, text width=3.5cm, align=left, font=\tiny] {\texttt{all\_gather} for \texttt{Shard}$\rightarrow$\texttt{Replicate} inputs; \texttt{reduce\_scatter} for \texttt{Partial} outputs---minimal set only};
\draw[arrow, color=tpblue] (user) -- (check);
\draw[arrow, color=tppurple] (check) -- (reshard);
\draw[arrow, color=tpgreen] (reshard) -- (local);
\draw[arrow, color=tpblue] (local) -- (out);
\end{tikzpicture}
\caption{Sharding resolution: the sharding-aware layer between the user op and the native kernels. Placements are checked, the minimal collectives are inserted, and each shard then runs through the ordinary backend dispatcher.}
\label{fig:dtensor}
\end{figure}

\paragraph{Functional collectives with autograd.} Each collective is exposed twice: a plain function and an autograd-aware variant built as a differentiable \texttt{Function}, so distributed operations compose inside the ordinary DAG. Communication overlaps computation through a lazy tensor wrapper:

\begin{lstlisting}[caption={Lazy collective tensors trigger the underlying wait only on first real use, letting the stream run ahead.}, label=lst:asynccoll]
class AsyncCollectiveTensor(tp.Tensor):
    """Tensor storage with a pending collective completion handle."""
    def __new__(cls, elem, work=None):
        if isinstance(elem, AsyncCollectiveTensor):
            return elem            # already lazy; keep the single handle
        return _attach_async(elem, work)
    def trigger_wait(self, timeout=None):
        ...                        # waits once, then demotes to plain Tensor
        self.__class__ = tp.Tensor
\end{lstlisting}

The C++ side maps each Python function to a \texttt{ProcessGroup} method (\texttt{allreduce}, \texttt{all\_gather\_into\_tensor}, \texttt{reduce\_scatter\_tensor}, \texttt{broadcast}), with a gloo backend for CPU clusters and MPI for HPC interconnects; a symmetric-memory utility exposes NVLink/CXL zero-copy paths for the v2 collective implementations.

% Figure: load() format dispatch
\begin{figure}[H]
\centering
\begin{tikzpicture}[every node/.style={font=\scriptsize\sffamily}]
\node[libbox, fill=blue!8, draw=tpblue, minimum width=3.0cm, minimum height=0.7cm] (load) {\texttt{load()}};
\node[draw, diamond, aspect=2, fill=yellow!18, draw=tporange, below=0.45cm of load, inner sep=1pt, font=\tiny] (q1) {magic = MEGA?};
\node[libbox, fill=green!8, draw=tpgreen, below left=0.45cm and -0.35cm of q1, minimum width=3.2cm, minimum height=0.7cm] (mega) {native reader (mmap, digests)};
\node[draw, diamond, aspect=2.4, fill=yellow!18, draw=tporange, below right=0.45cm and -0.35cm of q1, inner sep=1pt, font=\tiny] (q2) {safetensors?};
\node[libbox, fill=green!8, draw=tpgreen, below=0.45cm of q2, minimum width=3.2cm, minimum height=0.7cm] (st) {JSON-header reader};
\node[libbox, fill=green!8, draw=tpgreen, below right=0.4cm and 0.0cm of st, xshift=1.7cm, minimum width=3.2cm, minimum height=0.7cm] (pt) {allowlist unpickler};
\node[libbox, fill=gray!12, draw=gray, below=1.55cm of q1, minimum width=3.4cm, minimum height=0.7cm] (map) {resolve \texttt{map\_location}};
\node[libbox, fill=blue!8, draw=tpblue, below=0.45cm of map, minimum width=3.4cm, minimum height=0.7cm] (ten) {tensors on target devices};
\draw[arrow, color=tpblue] (load) -- (q1);
\draw[arrow, color=tpgreen] (q1.west) -- ++(-0.35,0) |- node[above, pos=0.24, font=\tiny]{yes} (mega);
\draw[arrow, color=tporange] (q1.east) -- ++(0.35,0) |- node[above, pos=0.24, font=\tiny]{no} (q2);
\draw[arrow, color=tpgreen] (q2) -- node[right, font=\tiny]{yes} (st);
\draw[arrow, color=tporange] (q2.east) -- ++(0.4,0) |- node[above, pos=0.2, font=\tiny]{zip} (pt);
\draw[arrow, color=gray] (mega) |- (map);
\draw[arrow, color=gray] (st) |- (map);
\draw[arrow, color=gray] (pt) |- (map);
\draw[arrow, color=tpblue] (map) -- (ten);
\end{tikzpicture}
\caption{Checkpoint loading is format dispatch followed by one placement pass. Every branch is data-only; no branch executes embedded code.}
\label{fig:loaddispatch}
\end{figure}



%======================================================================
\subsection{Serialization: the .mega Format and Safe Interop}
\label{subsec:eval-serialization}

Model persistence is a security surface as much as a storage problem. TensorPlay's native format is a binary container with a four-byte magic, tagged metadata records (integers, floats, strings, nested arrays), storage deduplication, per-payload CRC32 or SHA256 digests, and a default 4\,096-byte alignment so payloads can be \texttt{mmap}-ed zero-copy. Loading is allowlist-based end to end: no pickle, no embedded code, ever.

Interoperability follows the same discipline. \texttt{safetensors} files are read and written through their JSON header format. Legacy zip-based checkpoints from other frameworks are parsed by a dedicated unpickler whose globals are strictly allowlisted to tensor-reconstruction functions; storage types map to native dtypes through a fixed table, and anything outside the allowlist raises instead of executing. Location handling (\texttt{map\_location}) accepts a device string, a device-to-device mapping, or a callable receiving a storage stub. A hub utility completes the lifecycle: cached downloads under a configurable root, hash verification, resume support, and atomic rename-on-complete.



%======================================================================
\subsection{Support Cores: Threading, BLAS Bridges, RNG}
\label{subsec:eval-supportcores}

Three medium-sized cores sit below the kernels but above the OS, each readable in one sitting.

\begin{table}[H]
\centering
\caption{Support cores.}
\label{tab:supportcores}
\small
\begin{tabular}{llr}
\toprule
Core & Content & Lines \\
\midrule
Intra-op pool & worker threads, condition-variable queue, grain-size \texttt{parallel\_for} & 452 \\
oneDNN context & one \texttt{dnnl::engine} per process, threadpool bridge & 137 \\
MKL/cBLAS bridge & \texttt{cblas\_sgemm}/\texttt{sgemv} behind \texttt{USE\_MKL} & (kernel file) \\
CPU RNG & self-contained MT19937 engine + generator state & 196 + 108 \\
CUDA RNG & Philox counter management, capture-aware state & 331 \\
\bottomrule
\end{tabular}
\end{table}

\paragraph{Intra-op pool.} The pool is created once and reused: the thread count moves through a \texttt{NOT\_SET $\rightarrow$ value $\rightarrow$ CONSUMED} atomic state machine so the first \texttt{set\_num\_threads} freezes the pool size, and \texttt{in\_parallel\_region} is answered from \texttt{thread\_local} flags instead of a mutex + thread-id scan. Kernels enter via a header template with a grain size; the C bridge exposes the same pool to oneDNN's partitioning.

\begin{lstlisting}[caption={The whole pool is one class: spawn workers once, drain on shutdown.}, label=lst:pool]
class IntraopThreadPool {
 public:
  explicit IntraopThreadPool(int num_workers) {
    for (int i = 0; i < num_workers; ++i) {
      threads_.emplace_back([this] { worker_loop(); });
    }
  }
  IntraopThreadPool(const IntraopThreadPool&) = delete;
  IntraopThreadPool& operator=(const IntraopThreadPool&) = delete;
  ~IntraopThreadPool() {
    {
      std::unique_lock<std::mutex> lk(mutex_);
      running_ = false;
      condition_.notify_all();
    }
    for (auto& t : threads_) t.join();
  }
  /* ... */
\end{lstlisting}

A dedicated harness verifies both properties empirically: an 8-thread/1-thread ratio of 0.94$\times$ on a 1-element add (target $\le$1.5$\times$: oversubscription does not hurt tiny ops) and 0.99$\times$ on a 4M-element add (parallelism engages when there is work).

\paragraph{oneDNN and MKL.} The oneDNN context holds one engine per process and, critically, drives oneDNN \emph{outside} its own callbacks: primitive partitioning is bridged onto the intra-op pool so nested parallelism cannot deadlock. CPU kernels call oneDNN eltwise primitives directly for tanh/relu-family activations, while matmul goes to MKL via \texttt{cblas\_sgemm} behind the \texttt{USE\_MKL} guard. This is the whole vendor-acceleration story of TensorPlay: two narrow bridges, no vendor data structures in the object model.

\paragraph{RNG.} CPU generators wrap a self-contained MT19937 so the distribution layer consumes raw engine words (the C++ standard's \texttt{std::distributions} are implementation-defined and therefore unusable for reproducibility):

\begin{lstlisting}[caption={The CPU engine is a literal Mersenne Twister pod; distributions read raw words.}, label=lst:mt]
// std::mt19937 because the distribution layer must consume the raw engine
// output directly (std::distributions are implementation-defined).
struct mt19937_data_pod {
    uint64_t seed_;
    int left_;
    bool seeded_;
    /* ... 624-word state, boxed-normal scratch ... */
\end{lstlisting}

The CUDA generator manages Philox counter state with capture-aware reservation, as shown in Listing~\ref{lst:philox}.

%======================================================================
\subsection{Python Bindings: The Full Surface}
\label{subsec:eval-bindings}

\S\ref{sec:arch} and \S\ref{subsec:eval-microbench} follow the bridge hot path. Table~\ref{tab:bindings} shows all 24 binding units, 12.6\,k lines in total.

\begin{table}[H]
\centering
\caption{Binding units, sorted by size. The import unit defines module order; the tensor unit owns method dispatch.}
\label{tab:bindings}
\small
\begin{tabular}{lrllr}
\toprule
Unit & Lines & Role & Unit & Lines \\
\midrule
tensor methods & 3{,}466 & methods, magic ops & device & 369 \\
bridge & 1{,}703 & parse/check/wrap & size & 339 \\
symbolic ints & 1{,}221 & symint arithmetic & bridge header & 279 \\
module init & 722 & import order, exceptions & dispatch & 187 \\
compile API & 670 & stax entry points & CUDA graphs & 173 \\
futures & 600 & Future plumbing & distributed autograd & 151 \\
distributed & 569 & c10d bindings & autocast mode & 107 \\
autograd & 517 & backward/grad hooks & transforms & 88 \\
namespace ops & 436 & free functions & dtype & 110 \\
RPC & 408 & RPC bindings & & \\
\bottomrule
\end{tabular}
\end{table}

The units least discussed elsewhere are the async and symbolic families. The futures unit binds the \texttt{Future} type used by RPC and collective completion callbacks; the distributed triple (distributed + RPC + distributed-autograd, 1{,}128 lines combined) front-ends the \texttt{distributed} Python tree of Table~\ref{tab:pybreakdown}; the symbolic-int unit exposes the same arithmetic the C++ core implements, so \texttt{stax} shape functions can run on unbacked symints during tracing.

%======================================================================
\subsection{Three-Track Test Suite}
\label{subsec:eval-tests}

The test tree---145 Python files (2{,}069 test functions, 2{,}571 collected including parametrized cases) plus 2{,}078 lines of C++ tests (gloo process-group, Vulkan api/op suites)---is organized as three families with distinct assertions:

\begin{table}[H]
\centering
\caption{The three-track system. Suffixes are naming conventions; each track asserts a different property of the same op.}
\label{tab:testtracks}
\small
\begin{tabular}{llrr}
\toprule
Track & Assertion style & Files & Tests \\
\midrule
\texttt{alignment} & numeric tolerance vs.\ reference & 13 & 152 \\
\texttt{native} & dispatch reached a backend-native kernel, not composite & 22 & 226 \\
\texttt{parity} & exact behavioral parity incl.\ exception types & 4 & 51 \\
\midrule
Other 106 files & engine, compiler, frontend, e2e & 106 & --- \\
\bottomrule
\end{tabular}
\end{table}

\begin{itemize}[leftmargin=1.2em]
\item \textbf{Alignment}: forward/backward outputs agree with the reference implementation within tolerance on a shape grid---the ``is the math right'' track.
\item \textbf{Native}: the op dispatched to a backend-native kernel, not the \texttt{Composite} fallback---the ``is it actually implemented'' track; this is the runtime complement of the operator census in \S\ref{subsec:eval-opcount}.
\item \textbf{Parity}: values, dtypes, strides, and error behavior match the reference exactly, including exception types---the ``contract'' track.
\end{itemize}

A full single-process run reports \textbf{9 failed, 2{,}305 passed, 263 skipped} in 163\,s. The 263 skips are environmental, not structural: 162 require absent CUDA hardware (module-level \texttt{skipif} guards), the rest require optional packages that are not installed. The 9 failures are localized in two files---antialias/nearest-exact upsample modes (6) and einsum contraction-order edge cases (3)---and are reported as open items rather than hidden.

\begin{figure}[H]
\centering
\begin{tikzpicture}[every node/.style={font=\scriptsize}]
% horizontal stacked bar of test outcomes (widths precomputed: 9.6 * n/2577)
\def\totalw{9.6}
\def\wp{8.5867}
\def\ws{0.9797}
\def\wf{0.0335}
\fill[tpgreen!25] (0,0.55) rectangle (\wp,1.05);
\draw[tpgreen] (0,0.55) rectangle (\wp,1.05);
\fill[gray!25] (\wp,0.55) rectangle (\wp+\ws,1.05);
\draw[gray] (\wp,0.55) rectangle (\wp+\ws,1.05);
\fill[red!25] (\wp+\ws,0.55) rectangle (\wp+\ws+\wf,1.05);
\draw[red!70!black] (\wp+\ws,0.55) rectangle (\wp+\ws+\wf,1.05);
\node[tpgreen!70!black] at (\wp/2,0.8) {passed 2{,}305};
\node[gray] at (\wp+\ws/2,1.35) {skipped (env) 263};
\draw[gray] (\wp+\ws/2,1.15) -- (\wp+\ws/2,1.05);
\node[red!70!black, anchor=west] at (\wp+\ws+0.15,1.35) {failed 9};
\draw[red!70!black] (\wp+\ws+\wf,1.15) -- (\wp+\ws+\wf,1.05);
\node[anchor=east,gray] at (0,0.3) {outcome of one full-suite run:};
\node[anchor=west,gray] at (\totalw+0.2,0.3) {2{,}577 total};
\end{tikzpicture}
\qquad
\begin{tikzpicture}[every node/.style={font=\scriptsize}]
% three-track test counts
\def\cw{0.78}
\foreach \x/\name/\count/\fillc/\drawc in {
  0/{*}\_alignment/152/blue!8/tpblue,
  1/{*}\_native/226/orange!10/tporange,
  2/{*}\_parity/51/purple!8/tppurple}{
  \pgfmathsetmacro{\h}{\count/100}
  \pgfmathsetmacro{\lbl}{\h+0.45}
  \fill[\fillc] (\x*\cw,0) rectangle (\x*\cw+0.5*\cw,\h);
  \draw[\drawc] (\x*\cw,0) rectangle (\x*\cw+0.5*\cw,\h);
  \node at (\x*\cw+0.25*\cw,\lbl) {\count};
  \node at (\x*\cw+0.25*\cw,-0.28) {\name};
}
\draw[->,gray] (-0.1,0) -- (2.35,0);
\end{tikzpicture}
\caption{Left: outcome of the full-suite run (163\,s). Right: test counts per track; 429 tests across the three suffix families, 2{,}069 across the whole tree.}
\label{fig:testoutcomes}
\end{figure}

%======================================================================
\subsection{Benchmark Suite: Measured Throughput}
\label{subsec:eval-bench}

The benchmark directory holds 20 harnesses. Table~\ref{tab:benche2e} reports five representative runs verbatim (16-core x86, 3.5\,GHz class, MKL + oneDNN + NNPACK enabled, 8 threads; the reference column is the mainstream CPU framework printed by the same script, so both sides pay identical environment costs). Speedup $=$ ref\,ms\,/\,tp\,ms.

\begin{table}[H]
\centering
\caption{End-to-end and layer benchmarks, 8 threads.}
\label{tab:benche2e}
\small
\begin{tabular}{llrrl}
\toprule
Harness & Workload & tp & ref & speedup \\
\midrule
llama e2e & 0.6B llama, prefill 128 (CPU) & 519.6\,ms & 612.6\,ms & 1.18$\times$ \\
 & decode 1 tok & 499.0\,ms/tok & 612.8\,ms/tok & 1.23$\times$ \\
training step & qkv-proj fwd/bwd & 114.7\,ms & 185.9\,ms & 1.62$\times$ \\
 & mlp-up fwd/bwd & 318.3\,ms & 469.5\,ms & 1.48$\times$ \\
 & mlp-down fwd/bwd & 338.2\,ms & 450.0\,ms & 1.33$\times$ \\
 & decode-1 (batch 1) & 24.2\,ms & 72.1\,ms & 2.98$\times$ \\
 & small-net fwd/bwd & 4.14\,ms & 4.44\,ms & 1.07$\times$ \\
optimizer step & Adam, 104 tensors & 6.22\,ms & 34.55\,ms & 5.56$\times$ \\
 & SGD & 5.20\,ms & 6.84\,ms & 1.31$\times$ \\
 & geomean, 10 optimizers & --- & --- & 3.83$\times$ \\
custom-op dispatch & overhead above bare kernel & 1.73\,\textmu s & 18.9\,\textmu s & 10.9$\times$ \\
intra-op pool & 1-elem add 8t/1t ratio & 0.94 & 0.94 & parity \\
 & 4M add 8t/1t ratio & 0.99 & --- & engages \\
\bottomrule
\end{tabular}
\end{table}

\begin{figure}[H]
\centering
\begin{tikzpicture}[every node/.style={font=\scriptsize}]
% speedup histogram of the hot-ops benchmark (58 op/shape pairs)
\def\cw{1.18}
\foreach \i/\lbl/\count/\fillc/\drawc in {
  0/{$<$0.8}/6/red!12/red!70!black,
  1/{0.8--0.95}/5/orange!12/tporange,
  2/{0.95--1.05}/3/gray!15/gray,
  3/{1.05--1.5}/16/blue!8/tpblue,
  4/{1.5--3}/16/tpgreen!15/tpgreen,
  5/{$\ge$3}/12/tpgreen!35/tpgreen!70!black}{
  \pgfmathsetmacro{\h}{\count/5.5}
  \pgfmathsetmacro{\lblh}{\h+0.3}
  \fill[\fillc] (\i*\cw,0) rectangle (\i*\cw+0.92*\cw,\h);
  \draw[\drawc] (\i*\cw,0) rectangle (\i*\cw+0.92*\cw,\h);
  \node at (\i*\cw+0.46*\cw,\lblh) {\count};
  \node at (\i*\cw+0.46*\cw,-0.32) {\lbl};
}
\draw[->,gray] (-0.15,0) -- (6.75,0);
\draw[gray] (3.68,3.55) -- (3.68,3.8) -- (6.6,3.8);
\node[anchor=west,gray] at (6.75,3.8) {ahead};
\node[gray] at (3.4,3.55) {parity};
\draw[->,gray,dashed] (2.05,3.5) node[above,font=\tiny]{geomean 1.73$\times$} -- (4.3,0.4);
\end{tikzpicture}
\caption{Per-op CPU throughput vs.\ the reference framework across 58 op/shape pairs (8 threads). 6 pairs trail by $>$20\% (gelu-tanh small shapes, small-batch \texttt{bmm}, \texttt{cumsum}); 28 pairs lead by $\ge$1.5$\times$, led by \texttt{sqrt} 10.8$\times$ and \texttt{norm2} 9.9$\times$ where the reference's scalar loop loses to the vectorized kernel.}
\label{fig:benchhist}
\end{figure}

Three readings of Figure~\ref{fig:benchhist} matter for the paper's claims. First, the distribution is wide on \emph{both} sides: a readability-first framework is not uniformly slower; on elementwise and reduction families the hand-written vectorized kernels win (exp 2.4$\times$, gelu 2.3$\times$ at 1M; 4.8$\times$/4.3$\times$ at 8M), while fused multi-kernel paths (gelu-tanh), small-batch \texttt{bmm}, and \texttt{cumsum} trail. Second, the wins and losses cluster by kernel file, which is exactly the readability promise: the pointwise kernel file is the single file to read to explain both tails. Third, matmul is reference-competitive (1K$^3$ 0.94$\times$, tall-skinny 1.8$\times$) because both frameworks land in the same MKL \texttt{cblas\_sgemm}---the dispatcher adds the same sub-1\% sliver measured in \S\ref{subsec:eval-microbench}.

The benchmark tree also contains heavier end-to-end harnesses: a 1{,}489-line ResNet-18 script that runs a 10-epoch training race plus a strict state-dict transfer requiring \texttt{allclose} logits and identical top-1 predictions, and a CUDA-graph replay harness that requires a GPU. Their CUDA-gated parts skip with the same environmental guards as the test tree.

%======================================================================
\subsection{Operator Census and Dispatch Depth}
\label{subsec:eval-opcount}

The operator census uses three counts: dispatcher registrations (\texttt{m.impl}), schema declarations (479 distinct names across 2{,}809 records), and the Python-visible surface (275 functions + 184 methods). The three-way comparison is the only way to distinguish ``kernel exists but unreachable'' from ``Python reachable but no kernel''.

\begin{table}[H]
\centering
\caption{Operator census.}
\label{tab:opcount}
\small
\begin{tabular}{lll}
\toprule
Layer & TensorPlay & Vendor baseline \\
\midrule
Registered kernels (\texttt{m.impl} dedup) & 673 (CPU 671, CUDA 587) & 2000+ distinct ops \\
Schemas & 479 names / 618 active & 2000+ (with overloads) \\
Backward formulas & 187 formulas $\rightarrow$ 148 nodes & thousands \\
Python \texttt{\_C} surface & $\sim$550 & $\sim$1500 \\
Coverage focus & pointwise, reduction, matmul, conv/pool (24 conv ops), linalg 26, \texttt{foreach} $\sim$100 & full domain \\
\bottomrule
\end{tabular}
\end{table}

The delta is deliberate. TensorPlay covers the dense tensor loop that a convolution classifier or a transformer pre-training step exercises, plus the composite fallbacks that make the remainder runnable: tensor properties, histogram, unique, itertools, and integration each register one backend-neutral \texttt{TENSORPLAY\_LIBRARY\_IMPL(Composite,\dots)} block that serves every backend until a native kernel overrides it. Missing families are documented as such rather than hidden: dropout arrived via a native kernel; \texttt{cross\_entropy}/\texttt{ctc\_loss}, \texttt{rms\_norm}, \texttt{scatter\_reduce}, and long-tail sparse ops remain composite-only or absent, each classified as ``kernel exists, schema missing'' or ``three layers missing'' so the gap is a table entry, not a surprise.

%======================================================================
\subsection{DispatchKey: 13 vs.\ 60+}
\label{subsec:eval-dispatchkey}

\begin{table}[H]
\centering
\caption{\texttt{DispatchKey} enumeration. Sentinel \texttt{EndOfKeys=13} is spelled out so the table size is a compile-time constant.}
\label{tab:dispatchkey}
\small
\begin{tabular}{lll}
\toprule
Key & Value & Role \\
\midrule
\texttt{CPU} & 0 & Dense backend \\
\texttt{CUDA} & 1 & Dense backend \\
\texttt{Vulkan} & 2 & Dense backend, reserved \\
\texttt{AutogradCPU} & 3 & \texttt{toAutogradKey(CPU)} = CPU + 3 \\
\texttt{AutogradCUDA} & 4 & Autograd over CUDA \\
\texttt{AutogradVulkan} & 5 & Autograd over Vulkan \\
\texttt{AutocastCPU} & 6 & \texttt{toAutocastKey(CPU)} = CPU + 6 \\
\texttt{AutocastCUDA} & 7 & Autocast over CUDA \\
\texttt{AutocastVulkan} & 8 & Autocast over Vulkan \\
\texttt{VmapCPU} & 9 & \texttt{toVmapKey(CPU)} = CPU + 9, outranks autograd \\
\texttt{VmapCUDA} & 10 & Vmap over CUDA \\
\texttt{VmapVulkan} & 11 & Vmap over Vulkan \\
\texttt{Composite} & 12 & Backend-neutral fallback \\
\texttt{EndOfKeys} & 13 & Sentinel, table size \\
\bottomrule
\end{tabular}
\end{table}

The dispatch table is therefore \texttt{array<atomic<KernelFunction>,13>} with each slot relaxed-initialized to \texttt{nullptr} and published with \texttt{memory\_order\_release}. The key set is a \texttt{uint32\_t} bitset over the 13 keys; priority resolution scans the mask for the numerically largest bit, so priority is \texttt{Vmap $>$ Autocast $>$ Autograd $>$ backend}. Helpers clear one axis a time, used by redispatch below each layer.

\begin{lstlisting}[caption={The \texttt{Composite} fallthrough in the hot path: one extra load, only when a backend slot is empty.}, label=lst:composite]
KernelFunction getKernel(DispatchKey key) const noexcept {
  auto k = table_->kernels[dispatchKeyIndex(key)].load(acquire);
  if (!k && is_backend_key(key)) {
    constexpr auto kCompositeIdx = dispatchKeyIndex(Composite);
    k = table_->kernels[kCompositeIdx].load(acquire);   // one extra load
  }
  return k;
}
\end{lstlisting}

TensorPlay keeps a compact set of dispatch keys and leaves specialized axes outside the core table. There is no \texttt{Meta} key and no open registration API: registration is a closed \texttt{unordered\_map<string,DispatchTable>} guarded by one mutex, and the \texttt{TENSORPLAY\_LIBRARY\_IMPL} macro expands to a process-lifetime static constructor. The cost is that shape-only inference, sparse dispatch, and quantized dispatch have no dedicated fast path; they either share the \texttt{Composite} fallback or are explicitly unsupported.

%======================================================================
\subsection{Autograd Nodes: 148 Generated, Explicit DAG, No Tape}
\label{subsec:eval-autograd-nodes}

The generator parses the derivative-formula registry (345 raw records collapsing to 187 formulas after grouping) through an expression AST and emits 148 node structs. Each node stores saved forward tensors as \texttt{SavedVariable} where the formula needs them and synthesizes \texttt{apply} from the same AST that decided the save set, so the save list, constructor, and \texttt{apply} cannot drift. Handwritten nodes that formulas cannot express (concatenation, stacking, rolling) follow the same \texttt{SavedVariable} contract.

There is no tape. The graph is the transitive closure of each node's \texttt{next\_edges}, where every \texttt{Edge} carries only hints---a shape hint, a gradient dtype, a device hint---so backward can materialize a missing gradient without touching the forward tensor. \texttt{GraphTask} holds a dependency counter map built once per backward, an input-buffer map keyed by raw \texttt{Node*}, and the captured output vector for \texttt{grad}. The \texttt{ReadyQueue} is a max-heap ordered by the node's creation sequence number: it pops the largest sequence number first, approximating reverse creation order, while true topological order is enforced by the dependency counter. The per-node interpreter executes ten steps---grad-mode guard, capture bookkeeping, missing-grad materialization, pre-hooks, evaluate-node guard, \texttt{apply}, post-hooks, \texttt{sum\_to\_shape}, dtype cast, NaN probe---then propagates to parents under the task mutex.

The memory/performance trade-off is explicit: every edge holds a \texttt{shared\_ptr<Node>} and every input buffer holds a \texttt{vector<Tensor>}, so one \texttt{shared\_ptr} per node plus one allocation per distinct input slot. A production runtime uses \texttt{intrusive\_ptr} to shave the control block. TensorPlay keeps the standard-library pointer so students can follow ownership with ordinary tooling; the constant-factor cost is a teaching decision, recorded as such.

%======================================================================
\subsection{Correctness Guards}
\label{subsec:eval-guards}

Four guards turn silent wrong gradients or silent cross-device copies into loud failures or into shape-correct grads. Table~\ref{tab:guards} summarizes them; the remainder of the subsection shows the exact check.

\begin{table}[H]
\centering
\caption{Correctness guards.}
\label{tab:guards}
\small
\begin{tabular}{lll}
\toprule
Guard & Check & Failure mode \\
\midrule
Version & \texttt{unpack} compares current vs.\ saved version & \texttt{RuntimeError} with both versions \\
Shape & \texttt{sum\_to\_shape} reduces broadcast dims & grad reshaped to the edge hint, no crash \\
Dtype & gradient cast to the recorded forward dtype & grad cast in low-precision backward \\
Device & device disagreement at op entry & distinct \texttt{DeviceMismatchError} \\
Anomaly & NaN probe + creation-stack capture & forward stack + backward chain printed \\
\bottomrule
\end{tabular}
\end{table}

\subsubsection{Version guard: \texttt{SavedVariable}}

Forward tensors needed by backward are not stored bare. \texttt{save} copies the tensor's version counter into the saved variable; \texttt{unpack} reloads the current version and throws on disagreement:

\begin{lstlisting}[caption={Version mismatch is loud: the message spells both versions.}, label=lst:savedvar-eval]
TP_THROW(RuntimeError,
  "one of the variables needed for gradient computation has been modified "
  "by an inplace operation: [saved version: " + to_string(saved_version_) +
  "; current version: " + to_string(current) + "]");
\end{lstlisting}

The version counter is an \texttt{atomic<uint32\_t>} bumped on every in-place op (guarded by ``not in inference mode''), and shared eagerly between a base and its views, so that mutating a view bumps the base's counter as well. Generated nodes store each tensor member as a \texttt{SavedVariable} and unpack the top of \texttt{apply}; handwritten nodes do the same. When the graph need not be kept, the engine releases the saved variables to free memory early.

\subsubsection{Shape and dtype guard: \texttt{sum\_to\_shape}}

Broadcast during forward inflates a scalar grad to the output shape. Before the next edge consumes it, the engine reduces it:

\begin{lstlisting}[caption={Broadcast recovery: reduce the leading dims and the size-1 broadcast dims, then reshape.}, label=lst:sumtoshape-eval]
static Tensor sum_to_shape(const Tensor& grad, const vector<int64_t>& target){
  if(target.empty()) return ops::sum(grad);                 // scalar ()
  int64_t leading = (int64_t)grad.dim() - (int64_t)target.size();
  if(leading < 0){ /* reshape up when numel matches, else passthrough */ }
  vector<int64_t> reduce; for(int64_t i=0;i<leading;++i) reduce.push_back(i);
  for(int64_t i=leading;i<(int64_t)grad.dim();++i)
    if(target[i-leading]==1 && grad.size(i)!=1) reduce.push_back(i);
  Tensor cur = reduce.empty()? grad : ops::sum(grad,reduce,true);
  if(leading>0) cur = ops::reshape(cur,target); return cur;
}
\end{lstlisting}

The recorded shape comes from the edge's shape hint (with ``no hint'' distinguished from the scalar shape \texttt{}), and the dtype comes from the edge's dtype hint. A floating-point grad whose dtype disagrees with the hint is cast back---the autocast case: an \texttt{fp32} grad re-entering a low-precision backward. Both ops run under the per-node grad-mode guard, so second-order grads still build a graph when \texttt{create\_graph} is true.

\subsubsection{Device guard: \texttt{DeviceMismatchError}}

Inference tensors never acquire an autograd key: their version counter is created disabled, so the key computation yields the bare backend bit. Mismatched devices are not silently copied; they throw a distinct Python type:

\begin{lstlisting}[caption={A distinct exception for device mismatch, derived from the runtime-error base.}, label=lst:devicemismatch]
class P10_API DeviceMismatchError : public RuntimeError {
  using RuntimeError::RuntimeError;
};
\end{lstlisting}

The throw fires on device disagreement in element-wise and matmul entry points, and in reduction iterators. Because the type derives from \texttt{RuntimeError} it carries the source location and an optional C++ stack, but Python can catch it separately from other runtime errors.

\subsubsection{Anomaly guard: \texttt{AnomalyMode}}

Enabling anomaly mode sets two process-global flags visible to every thread: record-creation-stacks and check-for-NaN. On node creation the mode stores the current stack (C++ frames by default; the Python binding installs \texttt{traceback.format\_stack} under the GIL) and, if a backward is currently evaluating, records the evaluating node as parent. During backward the engine pushes the evaluating node onto a thread-local guard, so any node created inside that \texttt{apply} records its parent. On exception the engine prints the forward stack chain before rethrowing. The NaN probe is

\begin{lstlisting}[caption={The NaN probe runs under a disabled grad-mode guard so it never records history.}, label=lst:nanprobe-eval]
bool tensor_has_nan(const Tensor& t){
  if(!t.defined() || !isFloatingOrComplexType(t.dtype())) return false;
  return t.ne(t).any().item<bool>(); // ieee NaN != NaN via dispatch
}
\end{lstlisting}

and the chain printer walks the recorded parent pointers (each a \texttt{weak\_ptr<Node>}), printing ``Previous calculation was induced by\dots'' per parent until expiration, so a \texttt{create\_graph} failure shows the forward site that induced the failing backward node.

%======================================================================
\subsection{Overhead When Disabled}
\label{subsec:eval-overhead}

Readable frameworks are judged by what they cost when the developer is not looking. Three hot paths sit on every op: the profiler gate, the autograd-mode gate, and the dispatch-key computation. All three are designed to be predicted not-taken or to be a pair of branches.

\begin{table}[H]
\centering
\caption{Disabled overhead budget, single thread.}
\label{tab:overhead}
\small
\begin{tabular}{llll}
\toprule
Gate & Instruction & Cost (disabled) \\
\midrule
Profiler & one \texttt{acquire} load + not-taken branch & $\approx$ 2 cycles \\
\texttt{GradMode} & \texttt{thread\_local bool} load + not-taken branch & $\approx$ 1 cycle \\
\texttt{InferenceMode} & \texttt{thread\_local bool} load + not-taken branch & $\approx$ 1 cycle \\
Dispatcher & 2 branches (autocast + composite) & $\approx$ 2 branches \\
Autograd key & grad-mode $\land$ inference-mode $\land$ requires-grad & 3 loads, short-circuit \\
\bottomrule
\end{tabular}
\end{table}

\subsubsection{Profiler: 2 cycles}

The profiler gate is one \texttt{memory\_order\_acquire} load inlined to a \texttt{mov} + \texttt{cmp} + predicted not-taken branch on x86. Shape and site capture each add one further \texttt{acquire} load, but they sit strictly inside the gate's true-branch, so they add zero cycles when disabled. A CPU-only session never arms the GPU timers at all: the binding stores the timing flags with \texttt{release} before the first op, and the GPU timer path early-returns when disarmed. Measured baseline overhead is $\sim$0.6\,ns at 3.5\,GHz for one load; see \S\ref{sec:profiler} for the enabled costs (30\,ns CPU, 80--120\,ns with shapes, two \texttt{cudaEventRecord} at $\sim$5\,\textmu s).

\subsubsection{\texttt{GradMode} / \texttt{InferenceMode}: \texttt{thread\_local}}

Both modes are \texttt{thread\_local bool}s (\texttt{GradMode} defaults true, \texttt{InferenceMode} false). Every generated wrapper computes \texttt{requires\_grad} as \texttt{grad\_mode \&\& !inference\_mode \&\& any\_input\_requires\_grad}---three inlined loads with short-circuit. Worker threads honor the task's \texttt{create\_graph} via a guard that overwrites TLS for the duration of node evaluation, so the disabled cost is exactly the TLS load the engine would have paid anyway. Inference tensors additionally skip version-counter allocation.

\subsubsection{Dispatcher: 2 branches}

\begin{lstlisting}[caption={Autocast exclusion is one predicted branch; composite fallthrough is the second.}, label=lst:dispbranch]
if (is_autocast_key(key) && !autocast::is_enabled(key)) [[unlikely]]
  TP_THROW(NotImplementedError, ...);
auto kernel = handle.getKernel(key);   // slot load; composite if empty
if (!kernel) TP_THROW(NotImplementedError, ...);
\end{lstlisting}

The first branch is the autocast exclusion above; the second lives inside \texttt{getKernel}: one \texttt{acquire} load of the indexed slot plus, only when a backend slot is empty, one extra \texttt{acquire} load of the \texttt{Composite} slot. For any op with a native kernel the composite load is not taken. Codegen hoists \texttt{static const OperatorHandle} so the name lookup is resolved once per process.

Two build-time choices amortize cold cost. The Python extension is compiled as one unity-build translation unit (the 9{,}300-line generated bindings header plus binding units together), while \texttt{p10} deliberately opts out of unity: kernel translation units must remain separate for per-file static registration. The shim library \texttt{tp\_python} merely re-exports real symbols, keeping import time dominated by \texttt{dlopen} rather than Python overhead.

The explicit DAG's teaching trade-off is one \texttt{shared\_ptr} per node versus \texttt{intrusive\_ptr}. A production runtime uses intrusive reference counting to avoid a separate control block per node; TensorPlay keeps the standard control block so ownership is observable in \texttt{gdb} and \texttt{valgrind} without custom pretty-printers. Sustained throughput therefore sees a constant-factor overhead in backward node allocation that is independent of tensor size and invisible on large matmuls where kernel time dominates.

%======================================================================
\subsection{Micro-Benchmarks: Dispatch and Binding Overhead}
\label{subsec:eval-microbench}

The disabled-overhead table above decomposes the C++ hot path into loads and branches. This subsection measures what a Python caller actually pays per op end-to-end, and where the nanoseconds go. All measurements are taken on the reference host (16-core x86-64, 3.5\,GHz class, quiet, single thread, min-of-reps over 2{,}000 iterations after 50 warm-up calls).

\subsubsection{Per-op wall-clock cost vs.\ tensor size}

\begin{table}[H]
\centering
\caption{Per-op wall-clock from Python, steady state (ns/op, min of 5 reps $\times$ 2{,}000 iters). The [16] fast path skips full \texttt{TensorIterator} build; [1024] adds kernel+memory traffic. \texttt{requires\_grad} adds the autograd wrapper layer.}
\label{tab:microbench-perop}
\small
\begin{tabular}{lrr}
\toprule
Operation & ns/op & Size-relative dispatch share \\
\midrule
\texttt{add} \texttt{[1]}, no grad & 885--1\,490 & dispatch $\sim$1\% of total \\
\texttt{add} \texttt{[1]}, \texttt{requires\_grad} & 1\,930 & autograd layer +$\sim$500\,ns (node alloc + edges) \\
\texttt{sum} \texttt{[1]} & 1\,360 & \\
\texttt{add} \texttt{[16]} & 870 & below [1]: warm cache, no iterator build \\
\texttt{add} \texttt{[1024]} & 940 & kernel+TLB begin to appear \\
\bottomrule
\end{tabular}
\end{table}

The first result to internalize: the \emph{C++ dispatch machinery costs single-digit nanoseconds} (Table~\ref{tab:overhead}: one acquire load, two predicted branches, one atomic slot index), but a Python-issued op pays $\sim$0.9--1.5\,\textmu s wall-clock. The gap is the binding layer: CPython argument tuple unpacking, bridge type checks, GIL-held entry, and result wrapping through the wrap cache. This is the expected shape for any Python-fronted framework and is why the FASTCALL binding path exists: it removes generic pybind11 dispatch ($\sim$150\,ns per call of overload resolution, measured by the gap between the module function at $\sim$890\,ns and a bare Python \texttt{lambda} adding two \texttt{int}s at $\sim$36\,ns on the same host).

\subsubsection{Layer-by-layer decomposition}

\begin{table}[H]
\centering
\caption{Where a Python \texttt{tp.add(a,b)} spends its $\sim$890\,ns (prebound module function vs.\ tensor method).}
\label{tab:microbench-decomp}
\small
\begin{tabular}{llr}
\toprule
Layer & Mechanism & Cost \\
\midrule
Python interpreter & bytecode dispatch, \texttt{lambda} call & $\sim$36\,ns (baseline) \\
Attribute + call setup & \texttt{METH\_FASTCALL} arg tuple & $\sim$80--120\,ns \\
Bridge parse/check & argument parse + type checks & $\sim$400\,ns \\
Key computation & \texttt{key\_set} + priority resolution & $\sim$2\,ns \\
Slot lookup & one \texttt{acquire} load & $\sim$0.6\,ns \\
Autocast+composite branches & predicted not-taken & $\sim$1--2\,ns \\
Result wrap & wrap-cache hit & $\sim$150\,ns \\
\midrule
Method over module fn & tensor method vs.\ prebound function & $+\sim$80\,ns \\
\bottomrule
\end{tabular}
\end{table}

Two conclusions follow. First, \textbf{the dispatcher is not the bottleneck}; C++-side dispatch is $\sim$4\,ns of the $\sim$890\,ns total ($<0.5\%$), so the 13-slot array and the closed-key design cost effectively nothing the Python layer where they would be noticed. Second, \textbf{the binding dominates}, which is intrinsic to per-op Python control flow: a training step that batches 1{,}000 ops into a compiled \texttt{stax} program replaces 1{,}000 bridge crossings with one (\S\ref{sec:arch:stax}), eliminating $\sim$0.9\,ms of pure per-op overhead per step on this host.

\subsubsection{CPU vs.\ CUDA vs.\ Vulkan dispatch overhead}

The three backends share one dispatcher table (\S\ref{sec:redispatch}); only the terminal kernel differs. The measured CPU numbers above are for \texttt{CPU}=slot 0. CUDA adds a stream push plus (if profiling) \texttt{cudaEventRecord}; Vulkan adds a command-buffer submit. Backend deltas follow from the kernel-entry structure: a CUDA op pays the same $\sim$4\,ns dispatch plus the driver launch constant ($\sim$5--10\,\textmu s host-side, the same order as the profiler's \texttt{cudaEventRecord} pair); a Vulkan op pays dispatch plus \texttt{vkQueueSubmit} ($\sim$10--30\,\textmu s). For all three backends the dispatcher's own contribution is the same sub-1\% sliver, which is the design goal: backend choice changes the kernel, never the dispatch path.

%======================================================================
\subsection{What Is Trimmed}
\label{subsec:eval-trimmed}

Table~\ref{tab:trimmed} enumerates the axes that TensorPlay cuts and what takes their place. Each cut is a one-line predicate in code or a single \texttt{Composite} registration, never a silent semantic change.

\begin{table}[H]
\centering
\caption{Trimmed axes. Every entry is a deliberate omission with a fallback or a documented gap.}
\label{tab:trimmed}
\small
\begin{tabular}{lll}
\toprule
Axis & What ships & What is trimmed \\
\midrule
Dispatch keys & 13 keys, \texttt{Composite} fallback & no third-party registration, no runtime key creation \\
\texttt{Meta} dispatch & none & no shape-only kernel; shape inference lives in the compiler pillar \\
Sparse / quantized & minimal COO storage, no quantized CPU ops & no sparse or quantized dispatch key \\
In-place alias analysis & \texttt{check\_inplace} on view functions & no multi-version alias tracker \\
Compiler & vertical pointwise fusion only & no loop codegen, no buffer planning, no autotune \\
Distributed autograd & collectives via Python c10d layer & no distributed autograd, no RPC autograd \\
Forward AD & 15 fused forward kernels & no per-level \texttt{ForwardADLevel} \\
\bottomrule
\end{tabular}
\end{table}

\begin{itemize}[leftmargin=1.2em]
\item \textbf{No third-party keys.} Registration is closed; there is no \texttt{registerDispatchKey} and no per-operator stack of boxed kernels. The benefit is that the name lookup is one hash probe plus one atomic load, not a vector walk.
\item \textbf{No meta dispatch.} A production \texttt{Meta} key runs kernels that only compute output shapes and dtypes. TensorPlay never dispatches meta kernels; shape propagation for the compiler pillar is ordinary Python with a shape-prop pass, dead-code elimination, and constant folding. The dispatcher therefore has no \texttt{Meta} slot and no ``meta-tensor'' concept.
\item \textbf{No sparse/quantized axes.} Sparse layout is a minimal storage payload (COO with a CSR variant), not a dispatch axis. Quantized affine int8 lives as a Python FakeQuantize straight-through estimator on CPU, not as a dtype$\times$key dispatch---although the Vulkan backend ships real quantized shaders (\S\ref{subsec:eval-backends}). There is no \texttt{SparseCPU} or \texttt{QuantizedCPU} key.
\item \textbf{Teaching compiler.} Fusion is exactly one pattern---\texttt{mul}$\rightarrow$\texttt{add}---plus pointwise expression programs with CPU vectorization. The graph executor interprets sequentially with lifetime-based buffer recycling; there is no codegen to loops or to CUDA, no buffer reuse planning, and no performance model. It is a concept demonstration, not a production compiler.
\item \textbf{Composite fallbacks.} The trimmed axes are not cliffs: two dozen composite registrations (view/copy, histogram, unique, itertools, integration, tensor properties) serve as the neutral kernel when a backend slot is empty. Adding a native kernel replaces the fallback for that op alone.
\end{itemize}
